\documentclass[12pt]{article}

% --- Packages ---
\usepackage[utf8]{inputenc}
\usepackage[T1]{fontenc}
\usepackage[french]{babel}
\usepackage{amsmath, amssymb}
\usepackage{geometry}
\usepackage{tikz}
\usepackage[most]{tcolorbox}
\usepackage{lmodern}
\usepackage{pgfplots}
\pgfplotsset{compat=1.18}
\usepackage{mdframed}
\usepackage{bm}
\usepackage{xcolor}
\usepackage{sectsty}
\usepackage[utf8]{inputenc}
\usepackage[T1]{fontenc}
\usepackage[dvipsnames]{xcolor}
\usepackage{tkz-tab}
\usepackage{multido}

\setlength{\parindent}{0pt}
\sectionfont{\color{red}\Large\bfseries}     

% --- Mise en page ---
\geometry{margin=2cm}

% --- pointillés -----
\newcommand{\Pointilles}[2][3]{%
  \par\nobreak
  \noindent\rule{0pt}{1.5\baselineskip}% Provides a larger gap between the preceding paragraph and the dots
  \multido{}{#2}{\noindent\makebox[\linewidth]{\rule{0pt}{#1\baselineskip}\dotfill}\endgraf}% ... dotted lines ...
  \bigskip% Gap between dots and next paragraph
}

\begin{document}

% ======== TITRE DU COURS ========
\begin{center}
    {\LARGE \textbf{Fonctions polynômes de degré 2}}\\[4mm]
    {\large 1STMG - Lycée Jacques Decour}
\end{center}

\hrule
\vspace{0.4cm}

% ============================
%   PARTIE 1 : Définition
% ============================

\section*{Partie 1 : Généralités}

\subsection*{1. Exemples et contre-exemples}

\[
\left.
\begin{aligned}
P_1(x) &= 3x^2 - 5x + 2 \\
P_2(x) &= -x^2 + 7x - 9 \\
P_3(x) &= (4+x)(5-2x) \\
\end{aligned}
\right\}
\quad
\text{Ce sont des polynômes de degré : }\_\_\_\_\_
\]

\[
\left.
\begin{aligned}
P_4(x) &= 5x - 3
\end{aligned}
\right\}
\quad
\text{C'est un polynôme de degré : }\_\_\_\_\_
\]

\[
\left.
\begin{aligned}
P_5(x) &= 4x^4 + 3x -2
\end{aligned}
\right\}
\quad
\text{C'est un polynôme de degré : }\_\_\_\_\_
\]

\vspace{10pt}

\begin{mdframed}[
    linewidth=1pt,
    roundcorner=5pt,
    linecolor=ForestGreen
]
{\color{ForestGreen}\underline{Définition :}}
    \\On appelle fonction polynôme de degré 2 toute fonction $f$ définie sur $\mathbb{R}$ par une expression de la forme :
\[
    f(x) = ax^2 + bx + c
\]
où $a$, $b$ et $c$ sont des réels et $a \neq 0$.
\end{mdframed}

\underline{Remarques :}
\vspace{5pt}
\begin{itemize}
    \item Une fonction polynôme du second degré s'appelle également \_\_\_\_\_\_\_\_\_\_\_\_\_\_\_.
\end{itemize}

\subsection*{2. Représentation graphique}

La représentation graphique d’une fonction polynôme du second degré est une \_\_\_\_\_\_\_\_\_\_\_\_\_\_.

\vspace{1pt}

\begin{mdframed}[
    linewidth=1pt,
    roundcorner=5pt,
    linecolor=red
]
    {\color{red}\underline{Propriétés :}}
    \vspace{5pt}
    \\Soit $f$ une fonction polynôme de second degré telle que $f(x) = ax^2+b$.
    \begin{itemize}
        \item si $a > 0$, $f$ est d'abord décroissante puis croissante (positif = content)
        \item si $a < 0$, $f$ est d'abord croissante puis décroissante (négatif = triste)
    \end{itemize}
\end{mdframed}

\bigskip

\begin{center}
\begin{tabular}{cc}
% --- GRAPHIQUE a > 0 ---
\begin{tikzpicture}[scale=1]
\begin{axis}[
    width=8cm,
    height=8cm,
    axis lines=center,
    grid=both,
    xmin=-4, xmax=4,
    ymin=-4, ymax=4,
    xtick={-4,-3,-2,-1,0,1,2,3,4},
    ytick={-4,-3,-2,-1,0,1,2,3,4},
    samples=200,
    title={$a > 0$},
]
\addplot[domain=-4:4, thick] {0.5*x^2 - 1};
\end{axis}
\end{tikzpicture}
&
% --- GRAPHIQUE a < 0 ---
\begin{tikzpicture}[scale=1]
\begin{axis}[
    width=8cm,
    height=8cm,
    axis lines=center,
    grid=both,
    xmin=-4, xmax=4,
    ymin=-4, ymax=4,
    xtick={-4,-3,-2,-1,0,1,2,3,4},
    ytick={-4,-3,-2,-1,0,1,2,3,4},
    samples=200,
    title={$a < 0$},
]
\addplot[domain=-4:4, thick] {-0.5*x^2 + 1};
\end{axis}
\end{tikzpicture}
\end{tabular}
\end{center}

\subsection*{3. Racines d'un polynôme}

\begin{mdframed}[
    linewidth=1pt,
    roundcorner=5pt,
    linecolor=ForestGreen
]
{\color{ForestGreen}\underline{Définition :}}
    \vspace{5pt}
    \\On considère une fonction polynôme du second degré : $f(x) = ax^2 + bx + c$.
    \\On appelle \textbf{racines du polynôme} $ax^2 + bx + c$ les solutions de l'équation $f(x) = 0$.
\end{mdframed}

\begin{mdframed}
    \underline{Exemple :}
    \vspace{5pt}
    \\On considère le polynôme du second degré $f(x) = 2x^2-3x+1$
    \vspace{5pt}
    \\On remarque facilement que $f(1) = 0$ car $2 \times 1^2 -3 \times 1 +1 = 0$
    \vspace{5pt}
    \\Donc $1$ est une \_\_\_\_\_\_\_\_\_\_\_\_ du polynôme $2x^2-3x+1$.
\end{mdframed}

\begin{mdframed}[
    linewidth=1pt,
    roundcorner=5pt,
    linecolor=ForestGreen
]
{\color{ForestGreen}\underline{Interprétation graphique :}}
    \vspace{5pt}
    \\Les racines d'un polynôme sont les antécédents de $0$, c'est-à-dire \textbf{les points d'intersection} de la \textbf{parabole} avec \textbf{l'axe des abscisses}.
\end{mdframed}

\begin{center}
\begin{tikzpicture}
\begin{axis}[
    width=14cm,
    height=12cm,
    axis lines=center,
    xmin=-5, xmax=5,
    ymin=-4, ymax=7,
    grid=both,
    samples=300,
    xlabel={$x$},
    ylabel={$y$},
    ytick={-4,-3,-2,-1,0,1,2,3,4,5,6,7},
    legend style={
        at={(1.05,1)},
        anchor=north west,
        draw=none
    }
]

% f1(x) = 3x^2 - 2x + 3
\addplot[thick, red] {3*x^2 - 2*x + 3};
\addlegendentry{$3x^2 - 2x + 3$}

% f2(x) = -2x^2 + 3x
\addplot[thick, blue] {-2*x^2 + 3*x};
\addlegendentry{$-2x^2 + 3x$}

% f3(x) = 5 - x^2
\addplot[thick, green!60!black] {5 - x^2};
\addlegendentry{$5 - x^2$}

% f4(x) = x^2 / 4
\addplot[thick, purple] {x^2/4};
\addlegendentry{$\dfrac{x^2}{4}$}

\end{axis}
\end{tikzpicture}
\end{center}

Le polynôme $3x^2 - 2x + 3$ possède \_\_\_\_\_\_ racine(s), qui sont \_\_\_\_\_\_\_\_\_\_\_\_ .
\vspace{10pt}
\\Le polynôme $-2x^2 + 3x$ possède \_\_\_\_\_\_ racine(s), qui sont \_\_\_\_\_\_\_\_\_\_\_\_ .
\vspace{10pt}
\\Le polynôme $5 - x^2$ possède \_\_\_\_\_\_ racine(s), qui sont \_\_\_\_\_\_\_\_\_\_\_\_ .
\vspace{5pt}
\\Le polynôme $\dfrac{x^2}{4}$ possède \_\_\_\_\_\_ racine(s), qui sont \_\_\_\_\_\_\_\_\_\_\_\_ .

% ============================
%   PARTIE 2 : Polynômes de la forme ax^2+c
% ============================

\section*{Partie 2 : Fonction polynômes du second degré de la forme $\bm{f(x) = ax^2+c}$ ($\bm{b=0}$)}

\subsection*{1. Fonction du type $f(x) = ax^2$}

\underline{1\up{er} cas :} {\color{red}$\bm{a>0}$}
\begin{center}
    \begin{tikzpicture}
    \begin{axis}[
        width=10cm,
        height=6cm,
        axis lines=middle,
        xmin=-5, xmax=5,
        ymin=0, ymax=12,
        grid=both,
        samples=200,
        xlabel={$x$},
        ylabel={$y$},
        xtick={-5,-4,-3,-2,-1,0,1,2,3,4,5},
        ytick={0,2,4,6,8,10,12},
        legend pos=north west,
        legend cell align=left
    ]
    
    
    % a = 0.5
    \addplot[thick, blue] {0.3 * x^2};
    \addlegendentry{$y=0{,}3x^2$}
    
    % a = 1
    \addplot[thick, red] {1 * x^2};
    \addlegendentry{$y=1x^2$}
    
    % a = 2
    \addplot[thick, green!60!black] {3 * x^2};
    \addlegendentry{$y=3x^2$}
    
    \end{axis}
    \end{tikzpicture}
\end{center}

\begin{mdframed}[
    linewidth=1pt,
    roundcorner=5pt,
    linecolor=ForestGreen
]
    La parabole est \textbf{tournée vers le haut}.
    \vspace{5pt}
    \\Plus $a$ est grand, plus les branches de la parabole se rapprochent.
    \\Plus $a$ est petit, plus les branches s'écartent.
    \\Avec $\bm{a>0}$, la fonction $f$ ne prend que des valeurs \textbf{positives}.
\end{mdframed}

On en déduit le tableau de signe :
\[
\begin{tikzpicture}
\tkzTabInit{$x$/1 , $f(x)$/1}{$-\infty$, , $+\infty$}
\tkzTabLine{,,,,}
\end{tikzpicture}
\]

\vspace{5pt}

\underline{2\up{ème} cas :} {\color{red}$\bm{a<0}$}
\begin{center}
    \begin{tikzpicture}
    \begin{axis}[
        width=10cm,
        height=7cm,
        axis lines=middle,
        xmin=-5, xmax=5,
        ymin=-12, ymax=2,
        grid=both,
        samples=200,
        xlabel={$x$},
        ylabel={$y$},
        xtick={-5,-4,-3,-2,-1,0,1,2,3,4,5},
        ytick={-12,-10,-8,-6,-4,-2,0,2},
        legend pos=south east,
        legend cell align=left
    ]
    
    
    % a = 0.5
    \addplot[thick, blue] {-0.3 * x^2};
    \addlegendentry{$y=-0{,}3x^2$}
    
    % a = 1
    \addplot[thick, red] {-1 * x^2};
    \addlegendentry{$y=-1x^2$}
    
    % a = 2
    \addplot[thick, green!60!black] {-3 * x^2};
    \addlegendentry{$y=-3x^2$}
    
    \end{axis}
    \end{tikzpicture}
\end{center}

\begin{mdframed}[
    linewidth=1pt,
    roundcorner=5pt,
    linecolor=ForestGreen
]
    La parabole est \textbf{tournée vers le bas}.
    \vspace{5pt}
    \\On a le même raisonnement pour le rapprochement et l'éloignement des branches.
    \\Avec $\bm{a<0}$, la fonction $f$ ne prend que des valeurs \textbf{négatives}.
\end{mdframed}

On en déduit le tableau de signe :
\[
\begin{tikzpicture}
\tkzTabInit{$x$/1 , $f(x)$/1}{$-\infty$, , $+\infty$}
\tkzTabLine{,,,,}
\end{tikzpicture}
\]

\begin{mdframed}[
    linewidth=1pt,
    roundcorner=5pt,
    linecolor=red
]
    {\color{red}\underline{\textbf{À RETENIR :}}}
    \vspace{5pt}
    \begin{itemize}
        \item le \textbf{signe de $\bm{a}$} détermine \textbf{les variations} (l'orientation de la parabole) et \textbf{le signe} de $f(x) = ax^2$.
        \item la \textbf{valeur de $\bm{a}$} détermine « l'ouverture de la parabole ».
    \end{itemize}
\end{mdframed}

\subsection*{2. Fonction du type $f(x) = ax^2 + c$}

\begin{minipage}{.55\textwidth}%
\begin{tikzpicture}
\begin{axis}[
    width=\textwidth,
    height=9cm,
    axis lines=center,
    grid=both,
    xmin=-3, xmax=3,
    ymin=-2, ymax=5,
    samples=200,
    xlabel={$x$},
    ylabel={$y$},
    legend pos=north west,
    legend cell align=left
]

% c = -1
\addplot[thick, blue] {0.3*x^2 - 1};
\addlegendentry{$y = 0,3x^2 - 1$}

% c = 0
\addplot[thick, red] {x^2};
\addlegendentry{$y = x^2$}

% c = 2
\addplot[thick, green!60!black] {3*x^2 + 2};
\addlegendentry{$y = 3x^2 + 2$}

\end{axis}
\end{tikzpicture}

\begin{tikzpicture}
\begin{axis}[
    width=\textwidth,
    height=9cm,
    axis lines=center,
    grid=both,
    xmin=-3, xmax=3,
    ymin=-4, ymax=3,
    samples=200,
    xlabel={$x$},
    ylabel={$y$},
    legend pos=north west,
    legend cell align=left
]

% c = -1
\addplot[thick, blue] {-3*x^2 - 1};
\addlegendentry{$y = -3x^2 - 1$}

% c = 0
\addplot[thick, red] {-x^2};
\addlegendentry{$y = -x^2$}

% c = 2
\addplot[thick, green!60!black] {-0.3*x^2 + 2};
\addlegendentry{$y = -0,3x^2 + 2$}

\end{axis}
\end{tikzpicture}
\end{minipage}%
\hfill
\begin{minipage}{.45\textwidth}%
Le sommet de la parabole est le point : $\bm{S(0,c)}$.
\vspace{5pt}
\begin{itemize}
    \item Si $\bm{c > 0}$ alors le polynôme $ax^2 + c$ n'a pas de racines.
    \item Si $\bm{c = 0}$ alors le polynôme a une
    \vspace{5pt}
    \\unique racine : \_\_\_\_\_\_\_\_\_\_.
    \item Si $\bm{c < 0}$ alors le polynôme $ax^2 + c$ a deux racines : $x_1$ et $x_2$. \textbf{$\bm{f}$ change de signe en $\bm{x_1}$ et $\bm{x_2}$}. On en déduit le tableau de signes suivant :
\end{itemize}
\vspace{10pt}
\[
\begin{tikzpicture}
\tkzTabInit[lgt=2, espcl=2.4]{$x$/1 , $f(x)$/1}{$-\infty$, , $+\infty$}
\tkzTabLine{,,,,}
\end{tikzpicture}
\]
\end{minipage}%

\begin{mdframed}[
    linewidth=1pt,
    roundcorner=5pt,
    linecolor=blue
]
    \underline{\textbf{Exercices d'application :}} 
    \vspace{10pt}
    \\\underline{Exercices 1 :}
    \vspace{5pt}
    \\Tracer les fonctions suivantes sur le cahier.
    \begin{enumerate}
        \item $f(x) = 3x^2+1$
        \item $g(x) = \dfrac{1}{2}x^2 - 2$
        \item $h(x) = 9 - x^2$
        \item $k(x) = -2x^2 - 3$
    \end{enumerate}
    \underline{Exercices 2 :}
    \vspace{5pt}
    \\On a représenté ci-dessous deux fonctions polynômes du second degré du type $f(x) = ax^2 + c$. Dans les deux cas, déterminer $a$ et $c$ en s'aidant du graphique.
    
    \vspace{10pt}
    
    \begin{minipage}{.4\textwidth}%
    \begin{tikzpicture}
    \begin{axis}[
        width=\textwidth,
        height=8cm,
        axis lines=center,
        grid=both,
        xmin=-3, xmax=3,
        ymin=-5, ymax=6,
        samples=200,
        xlabel={$x$},
        xtick={-3,-2,-1,0,1,2,3},
        ytick={-5,-4,-3,-2,-1,0,1,2,3,4,5},
        ylabel={$y$},
    ]
    
    % c = 2
    \addplot[thick, green!60!black] {2*x^2 - 4};
    
    \end{axis}
    \end{tikzpicture}
    \end{minipage}%
    \hfill
    \begin{minipage}{.6\textwidth}%
    \Pointilles[1.5]{8}
    \vspace{5pt}
    \end{minipage}%

    \begin{minipage}{.4\textwidth}%
    \begin{tikzpicture}
    \begin{axis}[
        width=\textwidth,
        height=8cm,
        axis lines=center,
        grid=both,
        xmin=-3, xmax=3,
        ymin=-4, ymax=4,
        samples=200,
        xlabel={$x$},
        xtick={-3,-2,-1,0,1,2,3},
        ytick={-4,-3,-2,-1,0,1,2,3,4},
        ylabel={$y$},
    ]
    
    % c = 2
    \addplot[thick, red!60!black] {-x^2 +3};
    
    \end{axis}
    \end{tikzpicture}
    \end{minipage}%
    \hfill
    \begin{minipage}{.6\textwidth}%
    \Pointilles[1.5]{8}
    \vspace{5pt}
    \end{minipage}%
    
\end{mdframed}

\vspace{10pt}

% ============================
%   PARTIE 3 : Polynômes du second degré sous forme factorisée $$
% ============================

\section*{Partie 3 : Polynômes du second degré sous forme factorisée $\bm{f(x) = a(x-x_1)(x-x_2)}$}

\vspace{5pt}
\begin{mdframed}[
    linewidth=1pt,
    roundcorner=5pt,
    linecolor=ForestGreen
]
{\color{ForestGreen}\underline{Définition :}}
\vspace{5pt}
    \\On considère une fonction polycôme du second degré de la forme : $f(x)=ax^2+bx+c$.
    \vspace{5pt}
    \\Si l'équation $f(x) = 0$ a deux solutions, c'est-à-dire si le polynôme $ax^2+bx+c$ a deux racines $x_1$ et $x_2$, alors la fonction $f$ peut s'écrire sous forme factorisée : 
\[
    f(x) = a(x-x_1)(x-x_2)
\]
\end{mdframed}

\underline{Remarque :}
\vspace{5pt}
\\Dans le cas où il n'y a qu'une seule racine $x_1=x_2$, alors la factorisation s'écrit : $f(x) = a(x-x_1)^2$.

\vspace{10pt}

\begin{mdframed}[
    linewidth=1pt,
    roundcorner=5pt,
    linecolor=red
]
{\color{red}\underline{Propriétés :} Étude de signes de la fonction $f(x)=a(x-x_1)(x-x_2)$}
    \vspace{5pt}
    \\
    \begin{minipage}{.45\textwidth}%
        Si $a>0$
        % \vspace{5pt}
        \[
        \begin{tikzpicture}
        \tkzTabInit[lgt=1.5, espcl=1.5]{$x$/1 , $f(x)$/1}{$-\infty$, $x_1$, $x_2$, $+\infty$}
        \tkzTabLine{,,z,,z,}
        \end{tikzpicture}
        \]
    \end{minipage}%
    \hfill
    \begin{minipage}{.45\textwidth}%
        Si $a<0$
        % \vspace{5pt}
        \[
        \begin{tikzpicture}
        \tkzTabInit[lgt=1.5, espcl=1.5]{$x$/1 , $f(x)$/1}{$-\infty$, $x_1$, $x_2$, $+\infty$}
        \tkzTabLine{,,z,,z,}
        \end{tikzpicture}
        \]
    \end{minipage}%
    
    \vspace{10pt}
    
    {\color{red}La courbe de la fonction $f(x)=a(x-x_1)(x-x_2)$ est symétrique par rapport à la droite d'équation $x = \dfrac{x_1+x_2}{2}$}

    \begin{minipage}{.45\textwidth}%
        \begin{tikzpicture}
        \begin{axis}[
            axis lines=center,
            width=\textwidth,
            height=8cm,
            % grid=both,
            xmin=-6, xmax=3,
            ymin=-8, ymax=5,
            samples=200,
            xlabel={$x$},
            ylabel={$y$},
            title={\textbf{Cas $\bm{a>0}$}},
        ]
        
        % Fonction
        \addplot[thick, blue] {(x+4)*(x-1)};
        
        % Zéros
        \addplot[only marks,mark=*] coordinates {(-4,0) (1,0)};
        \addplot[only marks, mark=*, color=orange] coordinates {(-1.5,0)};
        \node[below] at (axis cs:-4.5,1) {$x_1$};
        \node[below] at (axis cs:0.5,1) {$x_2$};
        
        % Droite pointillée : x = (x1+x2)/2 = 0.5
        \addplot[densely dashed, thick, orange] coordinates {(-1.5,-8) (-1.5,10)};
        \node[above right, orange] at (axis cs:-3.2,0) {$\frac{x_1+x_2}{2}$};
        
        \end{axis}
        \end{tikzpicture}
        \vspace{5pt}
        \\$f$ admet un minimum en $\dfrac{x_1+x_2}{2}$
    \end{minipage}%
    \hfill
    \begin{minipage}{.45\textwidth}%
        \begin{tikzpicture}
        \begin{axis}[
            axis lines=center,
            width=\textwidth,
            height=8cm,
            % grid=both,
            xmin=-6, xmax=3,
            ymin=-5, ymax=8,
            samples=200,
            xlabel={$x$},
            ylabel={$y$},
            title={\textbf{Cas $\bm{a<0}$}},
        ]
        
        % Fonction
        \addplot[thick, blue] {-(x+4)*(x-1)};
        
        % Zéros
        \addplot[only marks,mark=*] coordinates {(-4,0) (1,0)};
        \addplot[only marks, mark=*, color=orange] coordinates {(-1.5,0)};
        \node[below] at (axis cs:-4.5,1) {$x_1$};
        \node[below] at (axis cs:0.5,1) {$x_2$};
        
        % Droite pointillée : x = (x1+x2)/2 = 0.5
        \addplot[densely dashed, thick, orange] coordinates {(-1.5,-5) (-1.5,8)};
        \node[above right, orange] at (axis cs:-3.2,0) {$\frac{x_1+x_2}{2}$};
        
        \end{axis}
        \end{tikzpicture}
        \vspace{5pt}
        \\$f$ admet un maximum en $\dfrac{x_1+x_2}{2}$
    \end{minipage}%
    
\end{mdframed}

\begin{mdframed}[
    linewidth=1pt,
    roundcorner=5pt,
    linecolor=blue,
]
    \vspace{5pt}
    \underline{\textbf{Exercices d'application :}} 
    \vspace{10pt}
    \\\underline{\textbf{Exercices 3 :}} Étude algébrique de fonctions du type $f(x)=a(x-x_1)(x-x_2)$
    \vspace{5pt}
    \begin{enumerate}
        \item Étudier le signe et les variations de la fonction $f(x) = -2(x-2)(x+3)$.
        \item On considère la fonction $g$ définie par : $g(x) = 3x^2-9x+6$
        \begin{enumerate}
            \item Conjecturer les racines du polynôme $3x^2-9x+6$.
            \item En déduire la forme factorisée de $g$.
            \item Étudier le signe de la fonction de $g$.
            \item Donner les coordonnées du sommet de la parabole qui représente $g$.
        \end{enumerate}
    \end{enumerate}
    \vspace{10pt}
    \underline{\textbf{Exercices 4 :}} Étude graphique de fonctions du type $f(x)=a(x-x_1)(x-x_2)$\\
    \vspace{5pt}
    \begin{minipage}{.45\textwidth}%
    Déterminer la forme factorisée de la fonction $h$ représentée ci-contre.
    \end{minipage}%
    \hfill
    \begin{minipage}{.5\textwidth}%
    \begin{center}
    \begin{tikzpicture}
    \begin{axis}[
        width=\textwidth,
        height=8cm,
        axis lines=center,
        grid=both,
        xmin=-4, xmax=3,
        ymin=-5, ymax=4,
        xlabel={$x$},
        xtick={-4,-3,-2,-1,1,2,3},
        ylabel={$y$},
        ytick={-5,-4,-3,-2,-1,1,2,3,4},
        samples=300,
    ]
    
    % Fonction f(x) = 2(x-1)(x+2)
    \addplot[thick, blue] {2*(x-1)*(x+2)};

    \node[below, blue] at (axis cs:1.7,2) {$\mathcal{C}_h$};
    
    \end{axis}
    \end{tikzpicture}
    \end{center}
    \end{minipage}%
    
\end{mdframed}

\end{document}
